% compendium/execution-trading-costs/execution-algorithms.tex
% Entries migrated on 2026-10-15 from the legacy execution-algorithms file (slice 12), owned by the
% section overview execution-trading-costs; dispositions and source checks in notes/migration.md,
% the blocks and the original table verbatim in notes/artefacts/market-making-lob/derivations.tex.
% def:automated-trading-architecture added on 2026-10-15 from the legacy production-considerations file
% (slice 17), also owned by execution-trading-costs; its blocks verbatim in
% notes/artefacts/model-governance-controls/derivations.tex.

\section{Execution algorithms, order exposure and portfolio trading}

\begin{qfdefinition}[Automated trading architecture: market data with real-time analytics, two-way venue connections, historical data with recorded messages, and models with an accuracy check]\label{def:automated-trading-architecture}
\qffield{Statement} Lehalle and Laruelle sketch the main components of a broker's automated trading service (Figure~3.1, page image): tick-by-tick and context databases, a calendar of expiries, expected news and half-days, modelling and model selection, a trading framework with a strategic layer, local trading tactics and listeners at each venue, synchronisation, a trading station, a trading database and TCA \citep[sec.~3.1, pp.~193--194]{alpha-laruelle}. The inputs of such tools are four (sec.~3.1.1, pp.~194--197):
\begin{itemize}
\item \emph{market data}: real-time feeds, consolidated across venues but quick to disaggregate, because it matters which venue currently offers a price. They hold trades, each with a price, a quantity and a timestamp that is ``not always precise enough'', and order books, or the sequence of their modifications, from which the book must be rebuilt on the fly. They are enriched with analytics: the spread, the side that initiated a trade, intraday volatility, an efficient spread and estimates of hidden quantities. Since raw data and analytics are needed together, ``the sooner they are computed, the better'', ideally on hardware as soon as raw data arrive (pp.~194--195);
\item \emph{connections to venues}, which are not natively synchronised with market data, usually being faster, and are two-way. One direction sends new orders, cancellations and modifications; the other returns acknowledgements, rejections of ill-formed messages, cancellations and partial fills. Incoming messages also yield analytics: dark pools publish only trades, so a dark buy order filled in full at a price shows resting sell quantity at that price or better (pp.~195--196);
\item \emph{historical data}: tick-by-tick or order-by-order market data; the analytics already computed in real time; recordings of the messages the firm's algorithms sent and received, to study past performance; and context such as news, corporate events, circuit breakers and venue outages (p.~196);
\item \emph{models}, each with three components. A formula takes parameters and variables: in the impact formula $\kappa\sigma(v/V)^\gamma$ the parameters $\kappa,\gamma$ may be estimated once a day and stored, while the variables $\sigma$, $V$ and $v$ are, in the book's words, ``known in real time''. Interpretation: $V$, the market volume during the trading, is available only as it is observed; volume still to come requires a forecast (see the decision-time item in \hyperref[def:trading-envelope]{trading envelope}). An estimation procedure computes the parameters from historical data. An accuracy computation qualifies the model and its parameters on the immediate past or in real time, so that a model can be deactivated when it is unlikely to be adequate in the current context (pp.~196--197).
\end{itemize}

\qffield{Conditions} The inputs feed the components of a trading algorithm: execution parameters, a risk-control layer that sets the trading envelope, and tactics such as smart order routers (\hyperref[def:trading-envelope]{trading envelope}). For latency, some autonomous tactics sit as close as possible to the venues they address, and the risk-control layer at equal distance from all its robots (sec.~3.1.2, p.~198). This describes a broker's execution service for fragmented European equities as of the 2018 printing; it is a design, not a standard or a legal requirement.

\qffield{Practice} Interpretation: the recorded messages and context are what later make it possible to replay an execution and explain its performance (\hyperref[meth:execution-decision-support]{decision support}, \hyperref[meth:transaction-cost-analysis]{transaction cost analysis}). The accuracy computation is a model-monitoring control: governance of models and of the systems that run them is in \hyperref[meth:automated-trading-operational-controls]{operational controls for automated trading}.

\qffield{Pitfalls} The legacy principle's ``the hard part is not one model; it is keeping the decision chain coherent under fragmented venues and live latency'' is a synthesis, not a statement in the source. Several legacy items are not in the passages read: holidays, auctions, index events and expected halt days in the calendar (the figure shows expiries, expected news and half-days); book imbalance and queue sizes among the fast analytics; and public book updates as part of the venue connection, where the source keeps the private order messages apart from the market-data feed and notes that the two are not synchronised. Venue names, data and latencies are as of the 2018 printing.

\qffield{See also} \hyperref[def:trading-envelope]{trading envelope}; \hyperref[meth:execution-decision-support]{decision support}; \hyperref[def:smart-order-routing]{smart order routing} (dark probes, latency); \hyperref[meth:automated-trading-operational-controls]{operational controls for automated trading} (owned by \texttt{model-governance-controls}).

\qffield{Further reading} The legacy blocks it replaces, verbatim, are in \nolinkurl{derivations.tex} in the folder \nolinkurl{notes/artefacts/model-governance-controls/}.

\qffield{Index} trading architecture; trading tools; market data; consolidated feed; real-time analytics; venue connection; acknowledgement; rejection; partial fill; dark-pool probe; historical data; recorded messages; context database; calendar; model parameters; model variables; estimation; model accuracy; model deactivation

\qffield{Status} Verification: partly checked 2026-10-15 (against the text layer of a local copy of \texttt{alpha-laruelle}, 2018 World Scientific printing, secs.~3.1--3.1.2, pp.~193--198, page image of p.~194 with Figure~3.1). Stability: draft. Currency: practitioner description as of the 2018 printing.
\end{qfdefinition}

\begin{qfdefinition}[Trading algorithm: execution parameters, a risk-control layer that sets a trading envelope, and tactics inside it]\label{def:trading-envelope}
\qffield{Statement} Lehalle and Laruelle describe a trading algorithm, instantiated from a class such as VWAP, TWAP, implementation shortfall, target close or liquidity seeking, as three components \citep[sec.~3.1.2, pp.~197--198, p.~198 by page image]{alpha-laruelle}:
\begin{itemize}
\item \emph{execution parameters}: side and quantity, the benchmark (VWAP, TWAP, percent of volume, entry price, close price), start and end times, bounds on the trading rate (maximum and minimum percent), aggressiveness, access to dark pools;
\item a \emph{risk-control layer} that computes the \emph{trading curve}, a minimum and maximum trading rate, from models of market impact, market risk and alpha and from ``expected market contexts'', balancing impact against market risk and its adverse selection;
\item \emph{trading robots} or liquidity-seeking tactics, such as a smart order router, that navigate inside the envelope set by the risk-control layer using the immediate state of liquidity; the risk-control layer assigns them short-term goals.
\end{itemize}
Because the analytics behind the curve are noisy, the curve is thick: a \emph{trading envelope}, inside which no trajectory is a priori more efficient than another for risk control. An algorithm trading faster than its curve stops posting; one trading too slowly trades aggressively, paying the spread \citep[sec.~3.4.1, pp.~231--232, page image of p.~231]{alpha-laruelle}. Envelopes are shown for a VWAP and an implementation-shortfall algorithm at several risk levels (Figures~3.13--3.14, pp.~232--233).

\qffield{Notation} Interpretation: the envelope bounds the cumulative executed quantity, or equivalently its rate, as a function of time, as plotted in Figure~3.2 (p.~199) and Figure~3.13 (p.~232). It applies to the parent order; child-order size, price and venue are the tactics' decisions.

\qffield{Conditions} How each algorithm class pursues its benchmark, from the passages checked:
\begin{itemize}
\item \emph{Benchmark against policy.} VWAP and TWAP name both a price statistic and a policy (\hyperref[def:execution-benchmark]{execution benchmarks}). Kissell's VWAP policy trades in proportion to the intraday volume curve and his TWAP policy evenly in time; both are static schedules fixed for the day; his POV policy trades a set fraction of market volume, so it reacts to realised volume but does not guarantee completion, and some brokers add a minimum rate to finish by the end time \citep[chap.~1, pp.~16--17]{s7-kissellscience}. Arrival-price and implementation-shortfall algorithms solve a cost--risk optimisation, the latter adding real-time adaptation to volumes and prices; their meaning differs across brokers, as does the scale of the risk-aversion parameter \citep[same chap., pp.~17--19]{s7-kissellscience}.
\item \emph{Information at decision time.} Model parameters, such as $\kappa$ and $\gamma$ in the impact term $\kappa\sigma(v/V)^\gamma$, are estimated about once a day from history; the variables $\sigma$, $V$ and $v$ are fed in real time \citep[sec.~3.1.1, p.~197]{alpha-laruelle}. Interpretation: $V$, the market volume during the trading, is known only as it occurs, so a schedule fixed at the start uses a forecast of it, and a policy that adapts uses only the volume observed so far.
\item \emph{Participation.} Interpretation: neither source checked says whether the market volume a participation rate $\rho$ applies to includes the order's own fills. If it does, the order trades $\rho/(1-\rho)$ of other traders' volume (25 percent at $\rho=20$ percent); if not, $\rho$ of it, which is $\rho/(1+\rho)$ of the total (16.7 percent). Kissell's participation-weighted price, the volume-weighted price of the first (shares traded)$/\rho$ shares, counts the order's own fills in the window (p.~108). The measurement window, the denominator and auction volume (often skipped by TWAP, p.~17) need stating.
\end{itemize}

\qffield{Practice} Benchmarks can be combined by constraints: minimum or maximum participation rates make an algorithm behave like PoV when volume is very high or low; a limit price and a ``would level'' make it behave like implementation shortfall; a smart order router makes it resemble a liquidity seeker. Such hard combinations are bang-bang controls, applied fully once a threshold is crossed; they can be gamed if the threshold is detected, or adversely selected, finishing too early because of the constraints, if the conditions are not chosen accurately. Building the logic on envelopes avoids this. Customising then has two steps: inject the desired risk profile and any volume or price constraints into the envelope, and plug in tactics suited to the investment style under defined market conditions, such as ``go into dark pools if volatility increases'' \citep[sec.~3.4.1, pp.~231--233, page images]{alpha-laruelle}. The equations for optimal envelopes are cited there to Bouchard et al.\ (2011), Lehalle (2012) and Guéant et al.\ (2012), not consulted.

\qffield{Pitfalls} Lehalle and Laruelle refer to Appendix~A.6 for the trading curve (p.~197), but A.6 is the Harris tick-size model; the scheduling appendix is A.7 (Interpretation, from the table of contents and p.~266). The caption of Figure~3.14 reads ``for a VWAP (left) or a (right) trading algorithm'', with the second name missing; the text calls the pair a VWAP and an IS (p.~232). The legacy notes said a too-fast algorithm ``pauses or posts passively''; the source says it stops posting. Their PoV ``capped by urgency'', VWAP ``would-level'' constraints and IS with ``SOR and dark-pool probes'' rearrange the source's examples, in which the would level turns an algorithm into IS and the SOR into a liquidity seeker. The ``minimum-impact curve'' among the legacy benchmark paths is not in the passages checked. The legacy ``dark probes'' are supported: orders sent to dark pools come back as probes of their hidden books, and a firm records the messages its algorithms send and receive, with market data and context, to study their performance \citep[sec.~3.1.1, p.~196]{alpha-laruelle}; the legacy decision loop's full list of observed events (quotes, trades, fills, rejects, cancellations, venue responses) is a paraphrase, not a list in the source. Interpretation: the envelope is a risk-control construct; following it says nothing about the cost against a benchmark measured afterwards.

\qffield{See also} \hyperref[mod:almgren-chriss-execution]{mean--variance execution} (a model for the trading curve); \hyperref[meth:algorithm-benchmark-choice]{choosing an algorithm's benchmark}; \hyperref[def:execution-benchmark]{execution benchmarks}; \hyperref[meth:execution-decision-support]{pre-trade and monitoring tools}; \hyperref[def:price-impact]{price impact} (the models the risk-control layer uses, developed by \texttt{price-impact-execution}).

\qffield{Further reading} The legacy blocks it replaces, verbatim: \nolinkurl{notes/artefacts/market-making-lob/derivations.tex}.

\qffield{Index} trading algorithm; algorithm layers; execution parameters; risk-control layer; trading curve; trading envelope; tactics; trading robots; smart order router; SOR; liquidity seeking; benchmark customisation; minimum participation; maximum participation; limit price; would level; bang-bang control; static schedule; adaptive algorithm; participation rate; POV; PoV; VWAP algorithm; TWAP algorithm; risk aversion parameter

\qffield{Status} Verification: partly checked 2026-10-15 (against the text layers of local copies of \texttt{alpha-laruelle}, 2018 World Scientific printing, secs.~3.1.1--3.1.2, pp.~196--199, and sec.~3.4.1, pp.~231--233, page images of pp.~198, 231 and~233; and \texttt{s7-kissellscience}, Academic Press printing with copyright 2014, chap.~1, pp.~16--19, and chap.~3, p.~108). Stability: draft. Currency: practitioner descriptions as of 2014 and 2018; broker definitions vary.
\end{qfdefinition}

\begin{qfmethod}[Decision support for trading algorithms: pre-trade analytics, monitoring indicators and what-if scenarios]\label{meth:execution-decision-support}
\qffield{Statement} Before launch, pre-trade analytics use historical data, models and parameters to estimate and compare expected cost components (bid--ask spread, usual traded volume, usual quantity on the book), expected trading profiles and expected context (news, volatility levels), broken down by sector, currency and country; on portfolios they check consistency and speed some lines up and slow others down from the expected balance of market risk and impact \citep[sec.~3.1.3, pp.~198--199]{alpha-laruelle}. During trading, monitoring indicators let a trader, who may watch several hundred algorithms, find where poor performance concentrates; performance indicators estimate the result if trading stopped now or continued under unchanged conditions, from a decomposition of the average price into immediate price, market moves and impact (pp.~200--201, page image of p.~201). What-if scenarios give real-time estimates of changed execution conditions, such as finishing now by taking the visible liquidity or converting an implementation-shortfall algorithm to percentage of volume (p.~201). Pre-trade, monitoring and post-trade tools should be compatible: for an indicator such as the bid--ask spread, pre-trade gives its expected average, monitoring its current value, updated expectation and likelihood, and post-trade its value during trading and its influence on performance, with breakdowns by sector, currency, country and side (p.~202). Monitoring indicators are built for human decisions, such as changing execution parameters on the fly, and can also be used by models to identify causes of poor performance (p.~200). Because users must share the indicators' values with the algorithms, the algorithms themselves publish them as \emph{indicator publishers}, publishing views on their own state in real time to users and companion algorithms (p.~202, page image).

\qffield{Conditions} Some what-if outcomes are simple formulas, others are estimates (p.~201). Pre-trade figures are forecasts made before the order. The average-price decomposition ``follows'' the scheduling appendix's (A.26) (p.~201), but (A.26) is an ex-ante objective, $\mathbb E(W)+\lambda\operatorname{Var}(W)=v^\star S_0+\kappa\sum_n\sigma_nv_n^{\gamma+1}/V_n^\gamma+\lambda\sum_nx_n^2\sigma_n^2$, labelled immediate buy, market impact and market risk \citep[app.~A, sec.~A.7.2, p.~279, page image]{alpha-laruelle}: its third term is a variance penalty, not a price move. Interpretation: the realised counterpart, from the book's change of variable (p.~278), is the identity $W=v^\star S_0+\sum_nx_n\sigma_n\xi_n+\kappa\sum_n\sigma_nv_n^{\gamma+1}/V_n^\gamma$, dividing by $v^\star$ for the average price; it holds exactly inside the model (checked by simulation), with $x_n$ the quantity remaining from slice $n$ on, in the book's timing. On real executions the impact term is a model value and the market-move term what is left, so the split is model-dependent. The refined version, adding planned price, unexpected market moves and trading efficiency, is given in words only (p.~201), so its additivity cannot be checked.

\qffield{Practice} Interpretation: pre-trade estimates belong with the expected cost and risk of a schedule (\hyperref[mod:almgren-chriss-execution]{mean--variance execution}); realised performance belongs with \hyperref[def:implementation-shortfall]{implementation shortfall}. The average-price decomposition is the post-trade side of \hyperref[meth:transaction-cost-analysis]{transaction cost analysis}; its accounting counterpart without a model is the split in \hyperref[def:transaction-costs]{transaction costs}.

\qffield{Pitfalls} The legacy pre-trade taxonomy added ``portfolio coupling'' (common end times, exposure constraints, synchronisation); the source mentions synchronisation for portfolios with common constraints, such as the same end time or constant exposure to a factor, after the pre-trade list (p.~200), not in it. The legacy what-if note listed ``slow down, switch benchmark, seek dark liquidity, or cap participation''; the source's examples are finishing now and switching from IS to PoV. The legacy attribution taxonomy (slice 13) defined the immediate price as ``the reference available when the order decision is made'' and market impact as ``price movement caused by the order itself''; in the source the immediate price is $S_0$, the price of buying everything at once without impact, and the impact term is the model's value, not an observed causal effect. Its pre-, live- and post-trade principle follows p.~202. The legacy process \emph{Indicator Publication} (slice 17) follows pp.~200--202, but its conclusion that publishing makes algorithms ``stateful agents \dots\ rather than black boxes'' is its own: the book calls trading algorithms ``stateless agents'' with introspection (p.~202, checked on the page image) while having them publish views on their own state. The two wordings are left as found; the point the source makes is introspection and shared indicators.

\qffield{See also} \hyperref[def:trading-envelope]{trading envelope}; \hyperref[def:execution-benchmark]{execution benchmarks}; \hyperref[def:implementation-shortfall]{implementation shortfall}.

\qffield{Further reading} The legacy blocks it replaces, verbatim: \nolinkurl{notes/artefacts/market-making-lob/derivations.tex}.

\qffield{Index} pre-trade analytics; expected cost; trading profile; monitoring indicators; performance indicators; what-if scenarios; average price decomposition; planned price; trading efficiency; portfolio synchronisation; post-trade analysis; immediate price; unexpected market moves; performance attribution

\qffield{Status} Verification: partly checked 2026-10-15 (against the text layer of a local copy of \texttt{alpha-laruelle}, 2018 World Scientific printing, sec.~3.1.3, pp.~198--202, page images of pp.~198 and~201, and app.~A, sec.~A.7.2, pp.~278--279, page image of p.~279, added in slice 13; the realised identity checked by simulation; pp.~200 and~202, page image of p.~202, for indicator publication, slice 17). Stability: draft. Currency: practitioner description as of the 2018 printing.
\end{qfmethod}

\begin{qfdefinition}[Order exposure decisions: how a large trader shows, splits and routes an order, and the reactive--proactive continuum]\label{def:order-exposure-decisions}
\qffield{Statement} Order submission strategy is the most important determinant of execution quality that traders control: market against limit orders, where to place limits, and when and how to resubmit. Large traders must also decide whether to look for the other side or wait for it, whether to show their full interest or hide it, whether to split the order over time, whether to use one broker or several to hide their total interest, and whether to trade in one market or many; display decisions are the most important a large buy-side trader makes \citep[chap.~``Buy-Side Traders'', p.~380]{alpha-harris}. \emph{Reactive} traders trade in response to opportunities others present; \emph{proactive} traders articulate their interests. Proactive traders expose orders to attract reactive ones, and two reactive traders never trade with each other (same chap., sec.~18.3, pp.~383--384).

\qffield{Conditions} Cost is the main reason to be reactive: proactive traders issue instructions before knowing whether an opportunity will arise, which is wasteful when opportunities are rare, but reactive traders risk that proactive ones act first (p.~383). The reactive--proactive continuum is the willingness to bear search costs; the aggressive--passive continuum is the eagerness to trade. They are similar, not the same: aggressive traders tend to be proactive, and traders who can search cheaply tend to be proactive (p.~384).

\qffield{Practice} Evasive strategies: brokers to preserve anonymity, several brokers so that none knows the full interest, anonymous electronic systems, careful choice of whom to show first (the art of block trading), showing only part of the order, undisplayed orders, order indications that carry no commitment, and waiting for another trader's firm order and taking it \citep[same chap., sec.~18.5.1, pp.~386--387]{alpha-harris}. Deceptive and offensive strategies are costly and may be illegal in some jurisdictions (pp.~387--389).

\qffield{Pitfalls} The legacy principle equated proactive execution with crossing the spread; Harris separates search (reactive--proactive) from eagerness (aggressive--passive), and a proactive trader may post a limit order. Its claim that reactive execution ``usually lowers leakage'' fits Harris's evasive strategy of not exposing until trading (p.~387), but he states it as a strategy, not a regularity. The legacy lever list's ``discretion range'', ``protection price'', ``pegging'' and ``minimum fill'' are not in the passages checked. Systems and venues named (Autex, POSIT, Island and others) are as of 2002.

\qffield{See also} \hyperref[def:block-trade]{block trades} (order exposure as one of the four problems); \hyperref[def:order-anticipation]{order anticipation} (the parasitic traders exposure attracts); \hyperref[def:trading-envelope]{trading envelope} (tactics inside an algorithm). Fragmentation and dark venues are developed by \texttt{market-structure-execution-quality}.

\qffield{Further reading} The legacy blocks it replaces, verbatim: \nolinkurl{notes/artefacts/market-making-lob/derivations.tex}.

\qffield{Index} order submission strategy; order exposure; display decision; hidden orders; undisplayed orders; order indications; order splitting; multiple brokers; reactive trader; proactive trader; aggressive; passive; search cost; evasive strategies; deceptive strategies; offensive strategies

\qffield{Status} Verification: partly checked 2026-10-15 (against the text layer of a local copy of \texttt{alpha-harris}, 2002 Oxford printing, chap.~18, pp.~380--389). Stability: draft. Currency: concepts timeless; systems and legal remarks as of 2002.
\end{qfdefinition}

\begin{qfdefinition}[Portfolio trading: program trades, the firm bid/offer market for packages, basket algorithms and index-fund trading]\label{def:portfolio-trading}
\qffield{Statement} A \emph{program trade} is the simultaneous submission of many orders; for statistics the NYSE and SEC counted trades of 15 or more coordinated transactions worth \$1 million or more, about 27 percent of NYSE volume, of which index arbitrage was about 9 percent \citep[chap.~``Index and Portfolio Markets'', p.~489]{alpha-harris}. In the \emph{firm bid/offer market}, also called the package, basket or portfolio trading market, package dealers quote firm prices for a whole portfolio from a characterisation of it, usually relative to its end-of-day value (for example the close minus 15 cents a share), and the list is revealed only after the close \citep[same chap., sec.~23.4.1, p.~490]{alpha-harris}. A \emph{basket algorithm} manages cost against total basket risk, using correlations, rather than each stock's own risk; constraints include cash balancing, self-financing and minimum and maximum participation \citep[chap.~1, pp.~18--19]{s7-kissellscience}.

\qffield{Conditions} The characterisation states whether the portfolio is an index portfolio and, if not, summary statistics: quantities, firm sizes, betas, price levels, average trading volumes, volatilities, primary listings and index membership (p.~490). Dealers price portfolios better than single stocks because informed traders are less likely to trade portfolios; clients pay them to assume the trading problem, and dealers may place parts with other clients (pp.~490--491). Index products (index futures, index options and shares in index funds) let traders trade index risk more cheaply than the components, for three reasons Harris gives: index dealers rarely face well-informed traders, because few have valuable insight into the whole market; index markets are very active, so counterparties are easy to find and dealers turn over inventory quickly; and a single transaction is arranged, cleared and settled instead of many \citep[same chap., p.~484 and sec.~23.4, p.~489]{alpha-harris}. Index prices generally lead the cash index, because index traders price only index risk (p.~489).

\qffield{Practice} Index funds rebalance a value-weighted portfolio only when the index list changes, and a price-weighted one also at splits \citep[same chap., sec.~23.2, p.~486]{alpha-harris}. They can improve returns by supplying rather than demanding liquidity, substituting non-index stocks for expensive components, rebalancing only with flows and dividends, and holding a subset of the index; each tactic raises tracking error, so many funds use them sparingly (p.~487). A basket algorithm may slow a stock despite good liquidity, or speed one despite poor liquidity, if that reduces residual basket risk (Kissell, pp.~18--19).

\qffield{Pitfalls} The legacy taxonomy treated basket trades and package trades as different; Harris uses package, basket and portfolio trading as names for one market. It explained tight index-basket quotes by diversified residual risk; Harris's reason is less adverse selection (Interpretation: diversification may also help, but is not in the passage). Its characterisation inputs added sector, currency, country, factor loadings and borrow constraints, and its index-rebalance controls named futures and ETFs as substitutes and timing around the effective window; Harris names non-index stocks, and the timing around reconstitution is in \hyperref[ev:index-reconstitution-impact]{index reconstitution}. The legacy principle \emph{Index Markets Trade Common Risk} (slice 18) gave four reasons for lower costs: index flow is less security-specific, more active, easier to net and easier to hedge with futures, ETFs, options or baskets. The first two match Harris's adverse-selection and activity reasons; ``easier to net'' and ``easier to hedge'' are not among his reasons, and his third, a single transaction to arrange and settle, is missing from the legacy list. ``Portfolio transition'' is named by Kissell (``transition management'', p.~16) but not defined in the passages checked. The NYSE figures are as of 2002.

\qffield{See also} \hyperref[def:block-trade]{block trades} (single-stock size); \hyperref[ev:index-reconstitution-impact]{index reconstitution} (price pressure from forced index trading); \hyperref[mech:etf-creation-redemption]{ETF creation and redemption} (owned by \texttt{etf-index-closing-execution}); \hyperref[def:trading-envelope]{trading envelope}; \hyperref[meth:execution-decision-support]{pre-trade analytics} (portfolio consistency checks).

\qffield{Further reading} The legacy blocks it replaces, verbatim: \nolinkurl{notes/artefacts/market-making-lob/derivations.tex}.

\qffield{Index} program trade; program trading; package trading; basket trading; portfolio trading; firm bid/offer market; package dealer; portfolio characterisation; basket algorithm; portfolio algorithm; index fund; tracking error; index rebalancing; transition management

\qffield{Status} Verification: partly checked 2026-10-15 (against the text layers of local copies of \texttt{alpha-harris}, 2002 Oxford printing, chap.~23, pp.~486--491, with pp.~484 and~489 added in slice 18, and \texttt{s7-kissellscience}, Academic Press printing with copyright 2014, chap.~1, pp.~16--19). Stability: draft. Currency: market figures as of 2002.
\end{qfdefinition}

\begin{qfmechanism}[Positive-feedback hedging: portfolio insurance sells into falls and buys into rises, and its trades reinforce the move]\label{mech:positive-feedback-hedging}
\qffield{Statement} Portfolio insurance is a dynamic trading strategy that replicates a portfolio plus a put option; the formula for how much to trade comes from the Black--Scholes model and depends on the portfolio's volatility, the protection period and the guaranteed minimum value. Insurers buy when prices rise and sell when they fall, the same effect as stop orders, which many used. It works when the money covered is small, but when a large fraction of the market is insured, sellers arrive together and prices can fall substantially; order anticipators sell at the first sign of a fall and buyers withdraw, so expected insurance selling can drive prices down even before it occurs \citep[chap.~``Bubbles, Crashes, and Circuit Breakers'', sec.~28.2.2, pp.~560--561]{alpha-harris}. Interpretation (the legacy proposition): if a rule holds $q(S)$ shares with $q'(S)>0$, a fall $dS<0$ gives $dq=q'(S)\,dS<0$, a sale; if each trade moves the price by $\lambda\,dq$ with $\lambda>0$, the induced move has the sign of the original one.

\qffield{Conditions} Interpretation: if the induced move persists until the rule trades again and triggers the same response, the moves form a geometric series with ratio $\lambda q'(S)$, which is dimensionless ($\lambda$ in currency per share, $q'$ in shares per unit of currency); the total is $dS/(1-\lambda q'(S))$ for $\lambda q'<1$, for example $1.67\,dS$ at $\lambda q'=0.4$, checked by iteration. Purely temporary impact that reverts before the next rebalance gives no such amplification. Harris notes that exchange-listed puts would have made the insured amount visible through open interest and reduced the uncertainty that kept buyers away (p.~561).

\qffield{Practice} Harris's rough sizing for 1987: about \$100 billion insured, a typical delta of about 0.4, so a 1 percent fall would require selling about 0.4 percent, roughly 10 million shares at \$40, against average daily volume of 160 million shares (box, p.~561). Insurers who raised the protected value as markets rose had to sell more on the way down (box, p.~560).

\qffield{Pitfalls} Interpretation: the size of a rebalancing trade depends on how much the hedge ratio changes, the option's gamma, not on its delta. In a Black--Scholes illustration with zero rates, 20 percent volatility and one year to expiry, a 1 percent fall requires selling about 1.0, 1.6 and 1.8 percent of the insured value for protection at 80, 90 and 100 percent of the current value. Harris's box multiplies the fall by the delta, so its 10 million shares is an order-of-magnitude guess on a different basis. The legacy proposition shows only the sign of the effect, not its size or stability.

\qffield{See also} \hyperref[mech:bubble-crash-dynamics]{bubbles and crashes} (the 1987 crash, owned by \texttt{execution-market-making}); \hyperref[def:price-impact]{price impact}; \hyperref[def:portfolio-trading]{portfolio trading}.

\qffield{Further reading} The legacy blocks it replaces, verbatim: \nolinkurl{notes/artefacts/market-making-lob/derivations.tex}.

\qffield{Index} portfolio insurance; dynamic hedging; synthetic put; positive feedback trading; stop orders; hedging feedback; 1987 crash; delta; gamma; insurance overhang

\qffield{Status} Verification: partly checked 2026-10-15 (against the text layer of a local copy of \texttt{alpha-harris}, 2002 Oxford printing, chap.~28, pp.~560--561, with boxes; the feedback series and the Black--Scholes illustration are own calculations). Stability: draft. Currency: mechanism timeless; the sizing is historical, for 1987.
\end{qfmechanism}
